% GENERATED FILE. Edit canonical lessons, then run npm run generate:books.
% canonical-source-hash: fnv1a64:846ce2ef791d401f
% canonical-lessons: GU-C423-nagar, GU-C423-jillo, GU-C423-gatar, GU-C423-durgandh, GU-C423-vasavvu

\chapter{Towns and Drains}
\label{ch:gu-towns-and-drains}
\hlchaptermodality{423}

\begin{chapteropening}
\textbf{By the end of this chapter:} \emph{I can talk about a town and its drains.}

Complete the last lesson of chapter 423: I can talk about a town and its drains.
\end{chapteropening}

\section[\texorpdfstring{\gu{નગર} (nagar)}{નગર (nagar)}]{\gu{નગર} (nagar) — a town; a city (formal)}

\label{lesson:GU-C423-nagar}



\begin{quote}
Before the new one: say the Gujarati for a sweet-maker, then the Gujarati for to sell well.
\end{quote}

\subsection*{You'll want to know: \gu{નગર}}
\textbf{\gu{નગર}} — \emph{nagar} — \textquotedblleft{}a town; a city (formal)\textquotedblright{}.

\begin{quote}
\gu{આ} \gu{નગર} \gu{નદીના} \gu{કિનારે} \gu{છે}.
\end{quote}

\begin{quote}
\emph{ā nagar nadīnā kināre che.}
\end{quote}

\begin{quote}
This town is on the bank of the river.
\end{quote}

A neuter noun.

\begin{cousinweb}
\textbf{\gu{નગર}} (\emph{nagar}) is a tatsama, Sanskrit \emph{nagara}, \textquotedblleft{}a town, a city\textquotedblright{} (Monier-Williams); it ends many place names.
\end{cousinweb}

\begin{grammarlens}[title={What you've built}]
One more word for getting around.
\end{grammarlens}

\subsection*{Your turn}
\begin{itemize}
\raggedright
  \item \emph{Say it:} \emph{nagar}
  \item \emph{Say it:} \emph{nagar}, once more
  \item \emph{Say it:} \emph{ā nagar nadīnā kināre che.} — This town is on the bank of the river.
  \item \emph{Recall:} say the Gujarati for a sweet-maker, then the Gujarati for to sell well, then say \emph{nagar} again
\end{itemize}

\begin{tcolorbox}[breakable,colback=teal!4,colframe=teal!35!black,title={Before you move on}]
What does \gu{નગર} mean? (\textquotedblleft{}A town; a city (formal)\textquotedblright{}.) Say it once more.
\end{tcolorbox}
\section[\texorpdfstring{\gu{જિલ્લો} (jillo)}{જિલ્લો (jillo)}]{\gu{જિલ્લો} (jillo) — a district (of a state)}

\label{lesson:GU-C423-jillo}



\begin{quote}
Before the new one: say the Gujarati for to sell well, then the Gujarati for a town.
\end{quote}

\subsection*{You'll want to know: \gu{જિલ્લો}}
\textbf{\gu{જિલ્લો}} — \emph{jillo} — \textquotedblleft{}a district (of a state)\textquotedblright{}.

\begin{quote}
\gu{અમારા} \gu{જિલ્લામાં} \gu{ઘણાં} \gu{ગામ} \gu{છે}.
\end{quote}

\begin{quote}
\emph{amārā jillāmā̃ ghaṇā̃ gām che.}
\end{quote}

\begin{quote}
There are many villages in our district.
\end{quote}

A masculine noun.

\begin{cousinweb}
\textbf{\gu{જિલ્લો}} (\emph{jillo}) is usually traced to Arabic \emph{ḍilʿ}, \textquotedblleft{}a side, a rib\textquotedblright{}, which came through Persian as a word for a division of land.
\end{cousinweb}

\begin{grammarlens}[title={What you've built}]
One more word for getting around.
\end{grammarlens}

\subsection*{Your turn}
\begin{itemize}
\raggedright
  \item \emph{Say it:} \emph{jillo}
  \item \emph{Say it:} \emph{jillo}, once more
  \item \emph{Say it:} \emph{amārā jillāmā̃ ghaṇā̃ gām che.} — There are many villages in our district.
  \item \emph{Recall:} say the Gujarati for to sell well, then the Gujarati for a town, then say \emph{jillo} again
\end{itemize}

\begin{tcolorbox}[breakable,colback=teal!4,colframe=teal!35!black,title={Before you move on}]
What does \gu{જિલ્લો} mean? (\textquotedblleft{}A district (of a state)\textquotedblright{}.) Say it once more.
\end{tcolorbox}
\section[\texorpdfstring{\gu{ગટર} (gaṭar)}{ગટર (gaṭar)}]{\gu{ગટર} (gaṭar) — a drain; a sewer}

\label{lesson:GU-C423-gatar}



\begin{quote}
Before the new one: say the Gujarati for a town, then the Gujarati for a district.
\end{quote}

\subsection*{You'll want to know: \gu{ગટર}}
\textbf{\gu{ગટર}} — \emph{gaṭar} — \textquotedblleft{}a drain; a sewer\textquotedblright{}.

\begin{quote}
\gu{વરસાદમાં} \gu{ગટર} \gu{છલકાઈ} \gu{ગઈ}.
\end{quote}

\begin{quote}
\emph{varsādmā̃ gaṭar chalkāī gaī.}
\end{quote}

\begin{quote}
The drain overflowed in the rain.
\end{quote}

A feminine noun.

\begin{cousinweb}
\textbf{\gu{ગટર}} (\emph{gaṭar}) is borrowed from English \emph{gutter}; in Gujarati a \gu{ગટર} is any drain or sewer.
\end{cousinweb}

\begin{grammarlens}[title={What you've built}]
One more word for getting around.
\end{grammarlens}

\subsection*{Your turn}
\begin{itemize}
\raggedright
  \item \emph{Say it:} \emph{gaṭar}
  \item \emph{Say it:} \emph{gaṭar}, once more
  \item \emph{Say it:} \emph{varsādmā̃ gaṭar chalkāī gaī.} — The drain overflowed in the rain.
  \item \emph{Recall:} say the Gujarati for a town, then the Gujarati for a district, then say \emph{gaṭar} again
\end{itemize}

\begin{tcolorbox}[breakable,colback=teal!4,colframe=teal!35!black,title={Before you move on}]
What does \gu{ગટર} mean? (\textquotedblleft{}A drain; a sewer\textquotedblright{}.) Say it once more.
\end{tcolorbox}
\section[\texorpdfstring{\gu{દુર્ગંધ} (durgandh)}{દુર્ગંધ (durgandh)}]{\gu{દુર્ગંધ} (durgandh) — a stink; a bad smell}

\label{lesson:GU-C423-durgandh}



\begin{quote}
Before the new one: say the Gujarati for a district, then the Gujarati for a sewer.
\end{quote}

\subsection*{You'll want to know: \gu{દુર્ગંધ}}
\textbf{\gu{દુર્ગંધ}} — \emph{durgandh} — \textquotedblleft{}a stink; a bad smell\textquotedblright{}.

\begin{quote}
\gu{ગટરમાંથી} \gu{દુર્ગંધ} \gu{આવે} \gu{છે}.
\end{quote}

\begin{quote}
\emph{gaṭarmā̃thī durgandh āve che.}
\end{quote}

\begin{quote}
A bad smell is coming from the drain.
\end{quote}

A feminine noun.

\begin{cousinweb}
\textbf{\gu{દુર્ગંધ}} (\emph{durgandh}) is a tatsama, Sanskrit \emph{durgandha}, \textquotedblleft{}a bad smell\textquotedblright{} (Monier-Williams), made of \emph{dur-}, \textquotedblleft{}bad\textquotedblright{}, and \emph{gandha}, \textquotedblleft{}a smell\textquotedblright{}.
\end{cousinweb}

\begin{grammarlens}[title={What you've built}]
One more word for getting around.
\end{grammarlens}

\subsection*{Your turn}
\begin{itemize}
\raggedright
  \item \emph{Say it:} \emph{durgandh}
  \item \emph{Say it:} \emph{durgandh}, once more
  \item \emph{Say it:} \emph{gaṭarmā̃thī durgandh āve che.} — A bad smell is coming from the drain.
  \item \emph{Recall:} say the Gujarati for a district, then the Gujarati for a sewer, then say \emph{durgandh} again
\end{itemize}

\begin{tcolorbox}[breakable,colback=teal!4,colframe=teal!35!black,title={Before you move on}]
What does \gu{દુર્ગંધ} mean? (\textquotedblleft{}A stink; a bad smell\textquotedblright{}.) Say it once more.
\end{tcolorbox}
\section[\texorpdfstring{\gu{વસાવવું} (vasāvvũ)}{વસાવવું (vasāvvũ)}]{\gu{વસાવવું} (vasāvvũ) — to set up a home; to found (a town)}

\label{lesson:GU-C423-vasavvu}



\begin{quote}
Before the new one: say the Gujarati for a sewer, then the Gujarati for a stink.
\end{quote}

\subsection*{You'll want to know: \gu{વસાવવું}}
\textbf{\gu{વસાવવું}} — \emph{vasāvvũ} — \textquotedblleft{}to set up a home; to found (a town)\textquotedblright{}.

\begin{quote}
\gu{દાદાએ} \gu{ગામમાં} \gu{નવું} \gu{ઘર} \gu{વસાવ્યું}.
\end{quote}

\begin{quote}
\emph{dādāe gāmmā̃ navũ ghar vasāvyũ.}
\end{quote}

\begin{quote}
Grandfather set up a new home in the village.
\end{quote}

Past, with \gu{તેણે} (he or she …): \gu{તેણે} \gu{નવું} \gu{ગામ} \gu{વસાવ્યું}, \emph{teṇe navũ gām vasāvyũ}.

\begin{cousinweb}
\textbf{\gu{વસાવવું}} (\emph{vasāvvũ}) is the causative partner of \gu{વસવું}, to settle, a tadbhava of Sanskrit \emph{vasati}, \textquotedblleft{}dwells\textquotedblright{} (CDIAL 11435): to make people or a home settle in a place.
\end{cousinweb}

\begin{grammarlens}[title={What you've built}]
One more word for getting around.
\end{grammarlens}

\subsection*{Your turn}
\begin{itemize}
\raggedright
  \item \emph{Say it:} \emph{vasāvvũ}
  \item \emph{Say it:} \emph{vasāvvũ}, once more
  \item \emph{Say it:} \emph{dādāe gāmmā̃ navũ ghar vasāvyũ.} — Grandfather set up a new home in the village.
  \item \emph{Recall:} say the Gujarati for a sewer, then the Gujarati for a stink, then say \emph{vasāvvũ} again
\end{itemize}

\begin{tcolorbox}[breakable,colback=teal!4,colframe=teal!35!black,title={Before you move on}]
What does \gu{વસાવવું} mean? (\textquotedblleft{}To set up a home; to found (a town)\textquotedblright{}.) Say it once more.
\end{tcolorbox}
